\chapter{Συζήτηση και Κριτική Αξιολόγηση}

\section{Εισαγωγή στη Συζήτηση των Ευρημάτων}
Η παρούσα διπλωματική εργασία σχεδιάστηκε με γνώμονα τη διερεύνηση της πρακτικής εφαρμοσιμότητας και της προστιθέμενης αξίας των τεχνολογιών Ψηφιακού Διδύμου (\textlatin{Digital Twin}) σε σχολικές υποδομές. Μέσω της πιλοτικής εφαρμογής στα Εκπαιδευτήρια Πλάτων, η έρευνα επιχείρησε να μεταβεί από τη θεωρητική σύλληψη των Έξυπνων Πανεπιστημιουπόλεων (\textlatin{Smart Campuses}) σε μια απτή, μετρήσιμη και τεχνοοικονομικά αξιολογήσιμη εγκατάσταση. Τα αποτελέσματα που παρουσιάστηκαν στο προηγούμενο κεφάλαιο υποστηρίζουν την κεντρική ερευνητική υπόθεση: η ενοποίηση υλικού χαμηλού κόστους (\textlatin{Edge IoT}) με λογισμικό ανοιχτού κώδικα (\textlatin{openHAB, InfluxDB}) δύναται να δημιουργήσει έναν κλειστό βρόχο ελέγχου ο οποίος βελτιώνει ταυτόχρονα το μαθησιακό περιβάλλον και την ενεργειακή αποδοτικότητα. Στο παρόν κεφάλαιο επιχειρείται η εις βάθος ερμηνεία αυτών των ευρημάτων, η σύγκρισή τους με την υφιστάμενη βιβλιογραφία, καθώς και η ανάλυση των περιορισμών και των προκλήσεων που αναδείχθηκαν κατά την επιχειρησιακή λειτουργία του συστήματος.

\section{Ερμηνεία των Αποτελεσμάτων Ποιότητας Εσωτερικού Περιβάλλοντος (\textlatin{IEQ}) και Ακουστικής Άνεσης}
Η Ποιότητα Εσωτερικού Περιβάλλοντος (\textlatin{Indoor Environmental Quality - IEQ}) αποτελεί ίσως τον κρισιμότερο παράγοντα για την εξασφάλιση της υγείας, της γνωστικής απόδοσης και της ψυχολογικής ευεξίας των μαθητών και των εκπαιδευτικών. Ιδιαίτερα η ακουστική άνεση, η οποία συχνά παραβλέπεται κατά τον σχεδιασμό σχολικών κτιρίων σε σχέση με τη θερμική ή την οπτική άνεση, τέθηκε στο επίκεντρο της λειτουργίας του Ψηφιακού Διδύμου.

\subsection{Η Ψυχολογική Προσέγγιση: Παθητική Αυτορρύθμιση Συμπεριφοράς}
Σύμφωνα με τις οδηγίες του Παγκόσμιου Οργανισμού Υγείας (\textlatin{WHO}), τα επίπεδα θορύβου βάθους (\textlatin{background noise}) σε μια άδεια σχολική αίθουσα δεν πρέπει να υπερβαίνουν τα 35 \textlatin{dBA}, ώστε να εξασφαλίζεται επαρκής λόγος σήματος προς θόρυβο (\textlatin{Signal-to-Noise Ratio, SNR $\geq +15$ dB}) για την καθαρότητα της ομιλίας \cite{WHO1999}, \cite{Kephalopoulos2015}. Κατά τη διάρκεια της εκπαιδευτικής διαδικασίας, ωστόσο, τα επίπεδα αυτά ποικίλλουν έντονα. Η αρχική περίοδος αναφοράς (\textlatin{baseline}), όπου το σύστημα λειτουργούσε αμιγώς ως παθητικός καταγραφέας (\textlatin{Digital Shadow}), ανέδειξε συχνές υπερβάσεις του δυναμικού κατωφλίου των 55 \textlatin{dBA}.

Η καινοτομία της παρέμβασης δεν εντοπίζεται απλώς στην καταγραφή, αλλά στη φύση της ανατροφοδότησης (\textlatin{actuation}). Η επιλογή να αποφευχθεί η χρήση ηχητικών συναγερμών (\textlatin{buzzers}) εντός της τάξης υπήρξε απολύτως συνειδητή και επιστημονικά τεκμηριωμένη. Η χρήση ενός επιπρόσθετου ηχητικού ερεθίσματος σε ένα ήδη θορυβώδες περιβάλλον θα προκαλούσε περαιτέρω διάσπαση προσοχής, γνωστική υπερφόρτωση και αυξημένο στρες τόσο στους μαθητές όσο και στον εκπαιδευτικό, διακόπτοντας βίαια τη ροή του μαθήματος.

Αντιθέτως, η ενεργοποίηση μιας διακριτικής οπτικής ένδειξης (\textlatin{quiet-lamp}) λειτούργησε ως ένας μηχανισμός "παθητικής αυτορρύθμισης" (\textlatin{passive behavioral nudging}). Όταν το σύστημα αντιλαμβανόταν συνεχή υπέρβαση του ορίου των 55 \textlatin{dBA} για ένα προκαθορισμένο χρονικό διάστημα (π.χ. 20 δευτερόλεπτα), άναβε τη λυχνία. Η οπτική αυτή ειδοποίηση επέτρεψε στους μαθητές να αντιληφθούν μόνοι τους την αύξηση της έντασης και να χαμηλώσουν τον τόνο της φωνής τους, ενώ ταυτόχρονα παρείχε στον εκπαιδευτικό ένα αντικειμενικό εργαλείο διαχείρισης της τάξης. Η σημαντική μείωση στον αριθμό και τη διάρκεια των υπερβάσεων του ορίου που αποτυπώνεται στο Σχήμα~\ref{fig:noise_results} κατά τη δεύτερη φάση λειτουργίας, υποστηρίζει ότι οι τεχνολογίες \textlatin{IoT} δεν χρησιμεύουν μόνο για τη μηχανολογική διαχείριση ενός κτιρίου, αλλά μπορούν να διαδραματίσουν ενεργό παιδαγωγικό ρόλο, βελτιώνοντας την ακουστική άνεση και μειώνοντας την κόπωση (\textlatin{vocal fatigue}) των διδασκόντων.

\section{Ενεργειακή Βελτιστοποίηση και Τεχνοοικονομική Σκοπιμότητα στη Διαχείριση Εγκαταστάσεων (\textlatin{Facility Management})}
Στον άξονα της ενεργειακής αποδοτικότητας, το Ψηφιακό Δίδυμο μετατράπηκε από ένα εργαλείο ευεξίας σε ένα αυστηρό σύστημα Διαχείρισης Εγκαταστάσεων (\textlatin{Facility Management}) και Προγνωστικής Συντήρησης (\textlatin{Predictive Maintenance}). Η μείωση της μέσης ημερήσιας κατανάλωσης κατά 19,7\% (από 223 σε 179 \textlatin{kWh}/ημέρα) αποτελεί ένα πρακτικά σημαντικό και εμπορικά κρίσιμο εύρημα.

\subsection{Η Αντιμετώπιση των Κρυφών Φορτίων (\textlatin{Phantom Loads})}
Τα σχολικά συγκροτήματα παρουσιάζουν εξαιρετικά ασυνεχή προφίλ κατανάλωσης. Παρατηρούνται ακραίες αιχμές ζήτησης κατά τις πρωινές ώρες λειτουργίας και σχεδόν μηδενική ζήτηση τα απογεύματα, τα Σαββατοκύριακα και τις περιόδους διακοπών. Παρά ταύτα, η ανάλυση των δεδομένων της περιόδου αναφοράς (πριν την εφαρμογή του συστήματος) ανέδειξε σημαντική κατανάλωση ενέργειας κατά τις ώρες αδράνειας. Αυτό οφείλεται κυρίως στα «κρυφά» φορτία (\textlatin{phantom loads}), δηλαδή σε συσκευές που παραμένουν σε λειτουργία ή σε κατάσταση αναμονής χωρίς ουσιαστικό λόγο.

Χαρακτηριστικό παράδειγμα αποτέλεσε ο εξοπλισμός του κυλικείου. Μέσω της τοποθέτησης των έξυπνων μετρητών \textlatin{Shelly 1PM} και των αισθητήρων κατάστασης θύρας (\textlatin{reed switches}), το σύστημα άρχισε να καταγράφει τη συχνότητα ανοίγματος των θυρών των επαγγελματικών ψυγείων και την αντίστοιχη άσκοπη λειτουργία του συμπιεστή. Η εγκαθίδρυση του κανόνα \textlatin{FridgeGuardian} στο \textlatin{openHAB} (ενεργοποίηση τοπικού βομβητή όταν η πόρτα παρέμενε ανοιχτή πέραν ενός χρονικού ορίου) περιόρισε αισθητά τις απώλειες ψύξης. Ιδιαίτερα σημαντική ήταν η \textlatin{Year-over-Year} (\textlatin{YoY}) σύγκριση του Νοεμβρίου, όπου καταγράφηκε μείωση της τάξης του 32,5\% (από 252 \textlatin{kWh}/ημέρα τον Νοέμβριο του 2024 σε 170 \textlatin{kWh}/ημέρα τον Νοέμβριο του 2025). Αντιθέτως, οι μήνες Σεπτέμβριος και Οκτώβριος 2025 εμφανίζουν οριακή αύξηση μετρούμενης κατανάλωσης σε σχέση με το 2024 (βλ. Παράρτημα Β, Πίνακας Β.2), γεγονός που αποδίδεται στη διεύρυνση του αισθητηριακού δικτύου και στην παράλληλη λειτουργία του τοπικού \textlatin{MVP} και της \textlatin{cloud} πλατφόρμας κατά τη μεταβατική περίοδο ωρίμανσης του συστήματος. Η αιχμή αυτή της εξοικονόμησης υποστηρίζει την αποτελεσματικότητα του αμφίδρομου ελέγχου (κλείσιμο διακοπτών, απενεργοποίηση φορτίων εκτός ωραρίου) κατά τους πλέον ενεργοβόρους χειμερινούς μήνες.

Η ερμηνεία της έντονης μείωσης του Νοεμβρίου απαιτεί διαχωρισμό των συνιστωσών φορτίου. Με βάση την ανάλυση γεγονότων της πλατφόρμας, το ισχυρότερο αποτύπωμα εντοπίζεται στον περιορισμό των \textlatin{phantom loads} και στη βελτίωση της λειτουργίας του ψυκτικού εξοπλισμού του κυλικείου (χρόνος ανοιχτής πόρτας, κύκλοι συμπιεστή). Η συμβολή φορτίων θέρμανσης/κλιματισμού (\textlatin{HVAC}) δεν μπορεί να ποσοτικοποιηθεί πλήρως στο παρόν \textlatin{dataset}, καθώς δεν υπάρχει ανεξάρτητη υπομέτρηση ανά υποσύστημα. Παρόλα αυτά, η συνέπεια των καταγραφών \textlatin{FridgeGuardian} και των χρονοσειρών ισχύος υποστηρίζει ότι η κύρια συνεισφορά του 32,5\% προέρχεται από τη μείωση λειτουργικών απωλειών και όχι από μεμονωμένη εποχική διακύμανση.

\subsection{Ανάλυση Κόστους Κύκλου Ζωής (\textlatin{LCC}) και \textlatin{ROI}}
Η τεχνοοικονομική μελέτη (\textlatin{ROI Calculator}) κατέδειξε περίοδο απόσβεσης (\textlatin{Payback Period}) 42,5 μηνών και 5ετές \textlatin{ROI} της τάξης του 41,02\%\footnote{Σημείωση: Η οπτικοποίηση της πλατφόρμας (Σχήμα~5.1) εμφανίζει τιμή 41{,}44\% λόγω ελαφρά διαφορετικής στρογγυλοποίησης της παραμέτρου ετήσιας εξοικονόμησης. Η τιμή αναφοράς της παρούσας ανάλυσης είναι 41{,}02\%.}. Τα νούμερα αυτά αποκτούν ακόμη μεγαλύτερη βαρύτητα εάν συγκριθούν με τα παραδοσιακά Συστήματα Διαχείρισης Κτιρίων (\textlatin{Building Management Systems - BMS}). Ένα τυπικό \textlatin{BMS} κλειστής αρχιτεκτονικής συνεπάγεται αυξημένο αρχικό κόστος κεφαλαίου (\textlatin{CAPEX}) για καλωδιώσεις, ιδιόκτητα πρωτόκολλα (\textlatin{proprietary hardware}) και υψηλότερα κόστη συντήρησης (\textlatin{OPEX}) με ετήσιες άδειες χρήσης λογισμικού.

Αντιθέτως, το προτεινόμενο σύστημα βασίστηκε στην προσέγγιση του "\textlatin{Brownfield Integration}", δηλαδή στην ενσωμάτωση του υφιστάμενου εξοπλισμού σε ένα σύγχρονο \textlatin{IoT} δίκτυο, χωρίς να απαιτείται ξήλωμα ή συνολική αντικατάσταση των υφιστάμενων υποδομών. Η προσέγγιση αυτή μειώνει το κόστος μετάβασης, ενισχύει τη διαλειτουργικότητα, συνάδει με αντίστοιχες ευρωπαϊκές μελέτες για ενεργειακή βελτιστοποίηση κτιρίων \cite{Apostolopoulos2022}, \cite{Fokaides2020}, και εναρμονίζεται πλήρως με τις απαιτήσεις της αναθεωρημένης Ευρωπαϊκής Οδηγίας για την Ενεργειακή Απόδοση Κτιρίων \cite{EUParliament2024}.

Ιδιαίτερη σημασία έχει και η εξελικτική μετάβαση από το αρχικό \textlatin{MVP}, το οποίο εξυπηρετούσε κυρίως την τοπική συλλογή δεδομένων και την επιβεβαίωση των κανόνων ελέγχου, προς την τελική \textlatin{Cloud} πλατφόρμα ανάλυσης. Η μεταφορά των δεικτών κόστους, των αναφορών \textlatin{ROI} και των προβολών εξοικονόμησης σε ένα περισσότερο φιλικό, διαδικτυακά προσβάσιμο περιβάλλον βελτίωσε ουσιαστικά τη χρηστικότητα του συστήματος για μη τεχνικούς ρόλους, όπως ο διευθυντής της σχολικής μονάδας, ο οποίος μπορούσε πλέον να παρακολουθεί συγκεντρωτικά την απόδοση της επένδυσης χωρίς να εξαρτάται από τεχνικά εργαλεία τοπικής διαχείρισης.

\section{Θεωρητική Ανάλυση: Η Μετάβαση από την Ψηφιακή Σκιά στο Ψηφιακό Δίδυμο}
Στον ακαδημαϊκό χώρο της Πληροφορικής και της Μηχανικής, παρατηρείται συχνά μια εσφαλμένη ταύτιση της απλής τηλεμετρίας \textlatin{IoT} ή των τρισδιάστατων σχεδίων \textlatin{BIM} με τον όρο "Ψηφιακό Δίδυμο". Σύμφωνα με την ταξινόμηση που διατυπώθηκε αρχικά στο βιομηχανικό πλαίσιο \cite{Kritzinger2018} και εφαρμόζεται ευρέως στον κτιριακό τομέα \cite{Fuller2020}, \cite{Grieves2017}, απαιτείται σαφής διαχωρισμός τριών βαθμίδων:
\begin{itemize}
\item \textbf{\textlatin{Digital Model}:} Η στατική ψηφιακή αναπαράσταση χωρίς αυτόματη ροή δεδομένων (π.χ. ένα σχέδιο \textlatin{CAD}).
\item \textbf{\textlatin{Digital Shadow}:} Η ψηφιακή αναπαράσταση που λαμβάνει αυτόματα δεδομένα από το φυσικό αντικείμενο (μονόδρομη ροή - \textlatin{Sense}), όπως τα περισσότερα \textlatin{IoT Dashboards}.
\item \textbf{\textlatin{Digital Twin}:} Το σύστημα όπου υφίσταται αμφίδρομη ροή (\textlatin{bidirectional data flow}). Τα δεδομένα του φυσικού χώρου ενημερώνουν τον ψηφιακό, και οι αλγοριθμικές αποφάσεις του ψηφιακού χώρου επιφέρουν άμεση φυσική μεταβολή (\textlatin{Actuation}).
\end{itemize}

Η παρούσα εργασία γεφυρώνει αυτό το ακαδημαϊκό κενό, τεκμηριώνοντας πρακτικά την εφαρμογή της τρίτης και ανώτατης βαθμίδας. Μέσω της αρχιτεκτονικής του \textlatin{openHAB}, το οποίο λειτούργησε ως σημασιολογικό ενδιάμεσο λογισμικό (\textlatin{semantic middleware}), υλοποιήθηκε ο πλήρης κύκλος \textit{\textlatin{Sense-Decide-Act}}. Ο μικροελεγκτής \textlatin{ESP32} ανιχνεύει τον θόρυβο (\textlatin{Sense}), αποστέλλει τα δεδομένα μέσω \textlatin{MQTT} στον κεντρικό διακομιστή \textlatin{Raspberry Pi}, η μηχανή κανόνων (\textlatin{Rules Engine}) αποφασίζει βάσει του ορίου των 55 \textlatin{dBA} και του σχολικού ωραρίου (\textlatin{Decide}), και τέλος, το σύστημα αποστέλλει εντολή ελέγχου στο ρελέ \textlatin{Sonoff} για να ανάψει η φυσική λάμπα εντός της τάξης (\textlatin{Act}). Η διπλή αυτή κατεύθυνση των δεδομένων υποστηρίζει τον χαρακτηρισμό της πλατφόρμας ως ένα πραγματικό Ψηφιακό Δίδυμο, διαχωρίζοντάς την από τα συμβατικά συστήματα τηλεμετρίας (\textlatin{Digital Shadows}).

\section{Προκλήσεις, Περιορισμοί και Επεκτασιμότητα (\textlatin{Scalability})}
Η εξαγωγή συμπερασμάτων από μια μελέτη περίπτωσης (\textlatin{case study}) σε πραγματικό περιβάλλον οφείλει να συνοδεύεται από κριτική εξέταση των περιορισμών της εφαρμογής και των προκλήσεων που ενδέχεται να προκύψουν κατά την κλιμάκωση (\textlatin{scalability}) του συστήματος σε άλλα σχολικά συγκροτήματα.

\subsection{Τεχνικοί Περιορισμοί και Εξάρτηση Δικτύου}
Ο πρώτος σημαντικός περιορισμός αφορά την αξιοπιστία της δικτυακής υποδομής. Εφόσον το σύνολο των κόμβων \textlatin{IoT} (\textlatin{Edge Nodes}) χρησιμοποιεί το ασύρματο τοπικό δίκτυο (\textlatin{Wi-Fi}) του σχολείου για την επικοινωνία με τον \textlatin{MQTT Broker}, τυχόν αστάθεια ή διακοπή λειτουργίας των \textlatin{Access Points} (\textlatin{AP}) οδηγεί σε προσωρινή απώλεια τηλεμετρίας. Για την αντιμετώπιση αυτού του κινδύνου σε μελλοντικές εγκαταστάσεις μεγάλων διαστάσεων (\textlatin{multi-building campuses}), προτείνεται η χρήση πρωτοκόλλων χαμηλής ισχύος και μεγάλης εμβέλειας, όπως το \textlatin{LoRaWAN} ή η δημιουργία απομονωμένων δικτύων \textlatin{Zigbee/Matter}, τα οποία δεν εξαρτώνται από το εμπορικό δίκτυο \textlatin{Wi-Fi} του ιδρύματος.

\subsection{Τοπική Επεξεργασία στην Παρυφή και Αποδοτικότητα Δικτύου}
Μια από τις σημαντικότερες τεχνικές προκλήσεις κατά την εισαγωγή αισθητήρων (ειδικά μικροφώνων) σε σχολικές αίθουσες είναι η διατήρηση χαμηλού δικτυακού φόρτου και η σταθερότητα του συστήματος σε πραγματικές συνθήκες λειτουργίας. Για τον λόγο αυτό, η αρχιτεκτονική του συστήματος σχεδιάστηκε με έμφαση στην τοπική επεξεργασία (\textlatin{edge processing}) αντί της αποστολής ακατέργαστων ροών στο \textlatin{cloud}. Η επιλογή αυτή μειώνει σημαντικά το απαιτούμενο εύρος ζώνης και αποτρέπει την υπερφόρτωση του δικτύου σε ώρες αιχμής.

Ο μικροελεγκτής \textlatin{ESP32} συλλέγει το ηχητικό σήμα μέσω του πρωτοκόλλου \textlatin{I2S} και υπολογίζει τοπικά τη Μέση Στάθμη Ηχητικής Πίεσης (\textlatin{Equivalent Continuous Sound Level - LAeq}) σε \textlatin{decibel} (\textlatin{dBA}) εντός πτητικής μνήμης (\textlatin{RAM}). Προς τον κεντρικό διακομιστή αποστέλλεται αποκλειστικά η συμπυκνωμένη αριθμητική τιμή (π.χ. 58.2 \textlatin{dBA}), γεγονός που περιορίζει τον όγκο τηλεμετρίας ανά κόμβο. Σε συνδυασμό με την ασφάλεια δικτύου (κρυπτογραφημένες ζεύξεις, έλεγχος πρόσβασης και απομόνωση υπηρεσιών), το σύστημα επιτυγχάνει ανθεκτική και αποδοτική λειτουργία σε επίπεδο \textlatin{campus}.

\subsection{Γενικευσιμότητα (\textlatin{Generalizability}) και Κουλτούρα Χρηστών}
Προκειμένου ένα τέτοιο σύστημα να επεκταθεί σε περιφερειακό επίπεδο (π.χ. σε όλα τα δημόσια σχολεία ενός Δήμου), πρέπει να υπερκεραστεί το εμπόδιο της εξοικείωσης των χρηστών. Το διδακτικό προσωπικό και η διοίκηση του σχολείου πρέπει να εκπαιδευτούν ώστε να αντιμετωπίζουν την πλατφόρμα όχι ως "εργαλείο παρακολούθησης της εργασίας τους", αλλά ως βοηθό για τη βελτίωση του μαθησιακού περιβάλλοντος. Η σχεδίαση της διεπαφής (\textlatin{3D Dashboard} σε \textlatin{React}) έλαβε υπόψη αυτή την ανάγκη, παρέχοντας ένα ιδιαίτερα φιλικό περιβάλλον χρήστη (\textlatin{UX}) που δεν απαιτεί γνώσεις προγραμματισμού ή μηχανικού αυτοματισμών, καθιστώντας την τεχνολογία προσιτή στον τελικό χρήστη.

\subsection{Τεχνικές Προκλήσεις και Μαθήματα από την Υλοποίηση}
Η πιλοτική εφαρμογή σε πραγματικό σχολικό περιβάλλον ανέδειξε πέντε κρίσιμα τεχνικά ζητήματα, η αντιμετώπιση των οποίων εμπλούτισε σημαντικά το αρχιτεκτονικό αποτύπωμα της πλατφόρμας.

Πρώτον, η αρχιτεκτονική συγχρονισμού δεδομένων προκαλούσε καθυστέρηση άνω των 500ms στην ενημέρωση του \textlatin{frontend}, με αποτέλεσμα μη αποδεκτή εμπειρία χρήστη για πραγματικού χρόνου εποπτεία. Η λύση υλοποιήθηκε μέσω \textlatin{WebSocket} συνδέσεων (\textlatin{Supabase Realtime}) με \textlatin{optimistic updates} και \textlatin{connection pooling}, μειώνοντας την καθυστέρηση σε επίπεδα κάτω των 100ms. Το κύριο δίδαγμα είναι ότι η υποδομή πραγματικού χρόνου πρέπει να σχεδιάζεται εξαρχής και όχι να προστίθεται εκ των υστέρων.

Δεύτερον, τα αρχικά τρισδιάστατα μοντέλα προκαλούσαν πτώση του ρυθμού καρέ (\textlatin{frame rate}) κάτω των 30fps σε μεσαίας κατηγορίας συσκευές, υποβαθμίζοντας την εμπειρία χρήστη. Η λύση περιελάμβανε υλοποίηση \textlatin{Level of Detail (LOD)}, \textlatin{lazy loading} τρισδιάστατων στοιχείων και βελτιστοποιήσεις \textlatin{React Three Fiber} (\textlatin{instanced meshes, object pooling}).

Τρίτον, η ετερογένεια πρωτοκόλλων (\textlatin{MQTT} για \textlatin{ESP32}, \textlatin{HTTP REST} για \textlatin{Shelly 1PM}, \textlatin{GPIO} για αισθητήρες επαφής) δημιουργούσε πολυπλοκότητα ολοκλήρωσης. Η υιοθέτηση του \textlatin{openHAB} ως ενιαίου \textlatin{semantic middleware} επέλυσε το ζήτημα, παρέχοντας κοινό \textlatin{REST API} για όλες τις συσκευές ανεξαρτήτως πρωτοκόλλου.

Τέταρτον, οι αισθητήρες θερμοκρασίας και υγρασίας εμφάνισαν σταδιακή απόκλιση (\textlatin{drift}) με την πάροδο του χρόνου, υποβαθμίζοντας αθόρυβα την ακρίβεια των δεδομένων. Η λύση περιελάμβανε αυτοματοποιημένες ρουτίνες βαθμονόμησης και αγωγούς επαλήθευσης δεδομένων (\textlatin{data validation pipelines}) με ανίχνευση εκτός-ορίων τιμών.

Πέμπτον, ο σχεδιασμός απομόνωσης δεδομένων ανά οργανισμό (\textlatin{multi-tenancy}) σε συνδυασμό με αποδεκτή απόδοση ερωτημάτων αποτέλεσε αρχιτεκτονική πρόκληση. Η αξιοποίηση των \textlatin{Row-Level Security (RLS)} πολιτικών της \textlatin{Supabase/PostgreSQL} με \textlatin{SECURITY DEFINER} βοηθητικές συναρτήσεις επέτρεψε την εφαρμογή ασφάλειας σε επίπεδο βάσης δεδομένων, απλοποιώντας σημαντικά τον κώδικα εφαρμογής.

Συνολικά, τα ανωτέρω τεχνικά μαθήματα επιβεβαιώνουν ότι η ανάπτυξη ενός πραγματικού Ψηφιακού Διδύμου σε εκπαιδευτικό περιβάλλον απαιτεί όχι μόνο θεωρητική γνώση αλλά και επαναληπτική μηχανική (\textlatin{iterative engineering}) σε συνθήκες πεδίου.

\subsection{Ηθικές Παράμετροι Εγκατάστασης}
Η εγκατάσταση αισθητηριακού δικτύου σε ενεργό εκπαιδευτικό χώρο με ανήλικους μαθητές εγείρει ζητήματα που απαιτούν ρητή αντιμετώπιση. Καταρχάς, το σύστημα δεν καταγράφει ηχητικό περιεχόμενο (φωνές, συνομιλίες): οι κόμβοι \textlatin{ESP32} υπολογίζουν αποκλειστικά τη στάθμη ηχητικής πίεσης (\textlatin{LAeq} σε \textlatin{dBA}) εντός πτητικής μνήμης και αποστέλλουν μόνο την αριθμητική τιμή στον \textlatin{broker}. Δεύτερον, η εγκατάσταση πραγματοποιήθηκε με πλήρη γνώση και έγγραφη συναίνεση της διοίκησης των Εκπαιδευτηρίων Πλάτων, η οποία ενημερώθηκε για τον σκοπό, τη φύση και τα όρια της συλλογής δεδομένων. Τρίτον, κανένα δεδομένο ταυτοποίησης προσώπου (όνομα, εικόνα, φωνή) δεν συλλέγεται ή αποθηκεύεται από το σύστημα, διασφαλίζοντας την πλήρη ανωνυμία των μαθητών και του εκπαιδευτικού προσωπικού.

\section{Ευθυγράμμιση με Ευρωπαϊκές Πολιτικές και το Όραμα των \textlatin{Smart Campuses}}
Τα ευρήματα της παρούσας μελέτης βρίσκονται σε ισχυρή εναρμόνιση με το μακροπρόθεσμο στρατηγικό πλαίσιο της Ευρωπαϊκής Ένωσης. Το πρόγραμμα \textit{\textlatin{DIGITAL EUROPE}} δίνει ρητή έμφαση στη δημιουργία Τοπικών Ψηφιακών Διδύμων (\textlatin{Local Digital Twins}) για αστικές και κοινοτικές υποδομές \cite{EuropeanCommission2023}. Η εφαρμογή αυτής της στρατηγικής σε εκπαιδευτικά κτίρια αποτελεί φυσική επέκταση του ευρωπαϊκού οράματος ψηφιακού μετασχηματισμού, ευθυγραμμιζόμενη παράλληλα με τις ενεργειακές απαιτήσεις της αναθεωρημένης Οδηγίας για την Ενεργειακή Απόδοση Κτιρίων \cite{EUParliament2024}.

Επιπρόσθετα, η εισαγωγή του Ψηφιακού Διδύμου στην εκπαιδευτική βαθμίδα δημιουργεί μετρήσιμες ευκαιρίες εφαρμογής STEM εντός του σχολικού χώρου. Η σχολική εγκατάσταση παύει να είναι απλώς ένα στατικό κέλυφος και μετατρέπεται σε ένα ``Ζωντανό Εργαστήριο'' (\textlatin{Living Lab}). Οι μαθητές έρχονται σε άμεση επαφή με ζωντανά δεδομένα (Big Data, IoT, περιβαλλοντικοί δείκτες), γεγονός που μπορεί να ενισχύσει τα προγράμματα εκπαίδευσης STEM (Επιστήμη, Τεχνολογία, Μηχανική, Μαθηματικά) και να καλλιεργήσει ενεργά την περιβαλλοντική τους συνείδηση. Η διπλωματική αυτή υποστηρίζει ότι τα Ψηφιακά Δίδυμα δεν αποτελούν προνόμιο αποκλειστικά της βαριάς βιομηχανίας (\textlatin{Industry 4.0}), αλλά μια ώριμη, προσβάσιμη τεχνολογία ικανή να μετασχηματίσει το σύγχρονο εκπαιδευτικό οικοσύστημα (\textlatin{Education 4.0}).

\section{Διαλειτουργικότητα (\textlatin{Interoperability}) και Ενσωμάτωση Υφιστάμενων Υποδομών}
Ένα από τα σημαντικότερα τεχνικά επιτεύγματα της παρούσας μελέτης είναι η απόδειξη της έννοιας της διαλειτουργικότητας (\textlatin{interoperability}) και της πρακτικής εφαρμογής του '\textlatin{Brownfield Integration}'. Στη βιομηχανία, η αναβάθμιση σε ένα ψηφιακό οικοσύστημα σπάνια συνεπάγεται την πλήρη αποξήλωση του παλαιού εξοπλισμού. Η πραγματική πρόκληση για έναν Μηχανικό είναι η ενοποίηση υφιστάμενων (\textlatin{legacy}) υποδομών με τις σύγχρονες τεχνολογίες Διαδικτύου των Πραγμάτων (\textlatin{IoT}).

Το σύστημα που αναπτύχθηκε υπερβαίνει κατά πολύ τον απλό 'Ψηφιακό Σκιασμό' (\textlatin{Digital Shadow}), ο οποίος περιορίζεται στη μονόδρομη ροή δεδομένων για την απλή οπτικοποίηση. Αντιθέτως, μεταβαίνει στο επίπεδο ενός πραγματικού Ψηφιακού Διδύμου (\textlatin{true Digital Twin}) μέσω της εγκαθίδρυσης αμφίδρομης επικοινωνίας και ελέγχου. Χαρακτηριστικότερο παράδειγμα αυτής της μετάβασης αποτελεί η επιτυχής ενσωμάτωση της υφιστάμενης υποδομής \textlatin{PLC (Programmable Logic Controller)} που διέθετε το σχολείο για τον έλεγχο των πυλών οχημάτων (\textlatin{vehicular gates}). Αντί να αντικατασταθούν τα βαρέα μοτέρ και οι ελεγκτές, τοποθετήθηκαν ρελέ ξηρής επαφής (\textlatin{dry-contact relays}) τα οποία επικοινωνούν μέσω πρωτοκόλλου \textlatin{MQTT}.

Η προσέγγιση αυτή επέτρεψε τον απομακρυσμένο έλεγχο των πυλών απευθείας από την τρισδιάστατη διεπαφή (\textlatin{3D Dashboard}) της πλατφόρμας. Σε συνδυασμό με τους \textlatin{MQTT-based} ενεργοποιητές (\textlatin{actuators}) για την ανατροφοδότηση του θορύβου μέσα στις σχολικές αίθουσες (έλεγχος οπτικής λυχνίας), η πλατφόρμα τεκμηριώνει την ικανότητά της να γεφυρώνει συστήματα διαφορετικών κατασκευαστών και γενεών. Η επιτυχία αυτού του εγχειρήματος αποδυναμώνει την αντίληψη ότι τα Έξυπνα Σχολεία (\textlatin{Smart Campuses}) απαιτούν εξαιρετικά υψηλά κόστη υλοποίησης και κλειστές αρχιτεκτονικές \textlatin{BMS}.

\section{Παιδαγωγική Αξία: Το Ψηφιακό Δίδυμο ως Εκπαιδευτικό Εργαλείο}
Ενώ τα περισσότερα Συστήματα Διαχείρισης Κτιρίων (\textlatin{BMS}) εξυπηρετούν αποκλειστικά τους διαχειριστές της εγκατάστασης (\textlatin{Facility Managers}), η εισαγωγή ενός Ψηφιακού Διδύμου σε ένα σχολικό συγκρότημα προσφέρει μια ανεκτίμητη, παράλληλη εκπαιδευτική διάσταση. Το ίδιο το κτίριο λειτουργεί ως «Ζωντανό Εργαστήριο» (εν. γλγ. \textlatin{Living Lab}) \cite{Li2024}, παρέχοντας στους μαθητές άμεση πρόσβαση σε πραγματικά δεδομένα περιβαλλοντικής παρακολούθησης στο πλαίσιο εφαρμογών \textlatin{STEM}.

\subsection{Εικονικά Εργαστήρια (\textlatin{Virtual Labs}) και Αναλυτική Μάθησης}
Οι αισθητήρες που έχουν εγκατασταθεί για τη λειτουργική διαχείριση παράγουν έναν τεράστιο όγκο ανοιχτών δεδομένων (\textlatin{big data}), τα οποία μπορούν να αξιοποιηθούν στη μαθησιακή διαδικασία. Για παράδειγμα, στο μάθημα της Φυσικής ή της Πληροφορικής, οι μαθητές μπορούν να αντλήσουν τις χρονοσειρές από τη βάση \textlatin{InfluxDB} και να μελετήσουν τη μεταβολή της θερμοκρασίας, να αναλύσουν τη διάδοση των ηχητικών κυμάτων (\textlatin{dBA}) μέσα στην τάξη, ή να υπολογίσουν την ηλεκτρική ισχύ (\textlatin{Watts}) που καταναλώνουν τα ψυγεία.

Αυτή η προσέγγιση προάγει την εκπαίδευση βάσει έργου (\textlatin{Project-Based Learning}). Η πλατφόρμα, μέσω των \textlatin{2D} και \textlatin{3D} δυναμικών κατόψεων (\textlatin{Floorplans}), επιτρέπει τη δημιουργία Εικονικών Εργαστηρίων (\textlatin{Virtual Labs}). Οι μαθητές οπτικοποιούν τις αόρατες μεταβλητές του χώρου τους, καλλιεργούν την περιβαλλοντική τους συνείδηση και κατανοούν πρακτικά τις έννοιες της Εξοικονόμησης Ενέργειας, βλέποντας σε πραγματικό χρόνο τις επιπτώσεις των πράξεών τους (π.χ. τι συμβαίνει στην κατανάλωση όταν αφήνουν μια πόρτα ανοιχτή) \cite{Li2024}.

\section{Προοπτικές Τεχνητής Νοημοσύνης (\textlatin{AI}) και Προγνωστικής Συντήρησης}
Με το θεμέλιο της αναπτυχθείσας πλατφόρμας να είναι πλέον λειτουργικό και σταθερό, διανοίγονται σημαντικές προοπτικές για τη μελλοντική της επέκταση. Ο όγκος των τηλεμετρικών δεδομένων που συγκεντρώνεται καθημερινά, καθιστά το σύστημα έτοιμο για την ενσωμάτωση αλγορίθμων Μηχανικής Μάθησης (\textlatin{Machine Learning - ML}) και αυτόνομων πρακτόρων (\textlatin{AI Agents}).

\subsection{Μετάβαση στο \textlatin{Cognitive Digital Twin}}
Η τρέχουσα υλοποίηση βασίζεται σε προκαθορισμένους, στατικούς κανόνες ελέγχου (\textlatin{rules-based logic}) στο ενδιάμεσο λογισμικό (\textlatin{openHAB}), όπως για παράδειγμα το όριο των 55 \textlatin{dBA}. Το επόμενο εξελικτικό βήμα είναι η μετάβαση σε ένα Γνωστικό Ψηφιακό Δίδυμο (\textlatin{Cognitive Digital Twin}) \cite{Semeraro2021}. Σε αυτό το μοντέλο, οι αλγόριθμοι Τεχνητής Νοημοσύνης θα αναλύουν ιστορικά δεδομένα για να εντοπίζουν δυναμικά τις ιδανικές συνθήκες.

Για παράδειγμα, η ανάλυση ανωμαλιών (\textlatin{Anomaly Detection}) θα μπορούσε να εξετάζει την κατανάλωση ισχύος των κινητήρων του ψυγείου σε συνάρτηση με την εσωτερική θερμοκρασία και τον αριθμό ανοιγμάτων της πόρτας. Εάν το σύστημα εντοπίσει ότι ο συμπιεστής λειτουργεί περισσότερο από το κανονικό για να διατηρήσει την ψύξη, μπορεί να σημάνει συναγερμό για επερχόμενη βλάβη (π.χ. απώλεια ψυκτικού υγρού) πολύ πριν το ψυγείο σταματήσει να λειτουργεί. Η Προγνωστική Συντήρηση (\textlatin{Predictive Maintenance}) αυτής της μορφής ελαχιστοποιεί τον χρόνο διακοπής λειτουργίας (\textlatin{downtime}) και μειώνει δραστικά τα κόστη συντήρησης του εξοπλισμού, επιβεβαιώνοντας τον ρόλο του Ψηφιακού Διδύμου ως ένα ιδιαίτερα ισχυρό εργαλείο για την επιχειρηματική διαχείριση σύγχρονων εγκαταστάσεων.

Σε επόμενο στάδιο, η ίδια υποδομή μπορεί να υποστηρίξει υβριδικά μοντέλα πρόβλεψης (\textlatin{time-series forecasting}) για ημερήσια κατανάλωση, ταξινόμηση συμβάντων λειτουργίας και αυτόματη βελτιστοποίηση \textlatin{setpoints} ανά ζώνη κτιρίου. Στο πλαίσιο αυτό, το \textlatin{AI Energy Forecasting} και οι τεχνικές \textlatin{Predictive Analytics} μπορούν να αποτελέσουν το ενδιάμεσο βήμα ανάμεσα στο σημερινό σύστημα υποστήριξης αποφάσεων και στο μελλοντικό \textlatin{Cognitive Digital Twin}. Παράλληλα, η ενσωμάτωση διεπαφών \textlatin{Voice Assistants} μπορεί να καταστήσει τις σύνθετες αναλύσεις περισσότερο άμεσα προσβάσιμες σε διοικητικούς χρήστες, μεταφέροντας τα αποτελέσματα της τεχνητής νοημοσύνης από το επίπεδο των τεχνικών \textlatin{dashboards} στο επίπεδο της καθημερινής λειτουργικής αξιοποίησης.

